\section{Data}
\label{sec:data}

\paragraph{Basketball.} Play-by-play for the 2019--20 through 2022--23 NBA seasons, regular season and playoffs, comes from a public Kaggle database \citep{walsh2026}. We reconstruct the five players on the floor for each team throughout every game from substitution events, inferring the players on the floor at the start of each period from the first events in which they appear. This yields \res{N} games with complete on-floor records, from which we compute every player's time on the floor and every pair's shared time. The same reconstruction produces \res{N} stints (intervals with no substitution) used to fit ridge-regularized adjusted plus-minus (RAPM). Points summed over a game's stints equal the final margin exactly in \res{x}\% of games and within two points in \res{x}\%. We did not trace the remaining discrepancies individually; they affect only the RAPM baseline, not the game-level training targets, which use final scores. The database does not cover 2023--24, so the scoring model is trained on games through 2022--23 and applied to rosters described by 2023--24 player features.

Player features come from season totals published by PBP Stats \citep{pbpstats} for 2018--19 through 2023--24 (regular season), combined with height, weight and listed position from the Kaggle database. Each player-season is described by 27 variables: minutes per game and share of games played; usage, true shooting and points per 36 minutes; the share of field-goal attempts at the rim, from short mid-range, from long mid-range, from the corners and from above the break; shrunken accuracy at the rim and from three; free-throw rate; assists per 36 minutes and the shares of assists leading to threes and to rim attempts; offensive and defensive rebounds, steals, blocks, turnovers and fouls per 36 minutes; height and weight; and indicators for guard, forward and centre. No on-court plus-minus quantity is included, so RAPM remains an independent benchmark. Features are standardized using players with at least 200 minutes.

\paragraph{Hockey.} Play-by-play, shift charts and team rosters for the 2021--22 through 2024--25 NHL seasons come from the league's public endpoints \citep{nhlapi}. The 4 Nations Face-Off began on 12 February 2025, and every game on or after that date is removed from all training data and from all 2024--25 features. Shift charts give each skater's time on ice second by second, from which we compute time shares, pairwise shared time and stints, yielding \res{N} games and \res{N} stints; goals summed over stints equal the final margin in \res{x}\% of games. A shot-quality model (logistic regression on distance, angle, shot type and rebound status, fit on \res{N} unblocked attempts) assigns an expected-goal value to every unblocked shot. Each player-season is described by 26 variables: time on ice per game and share of games played; shots on goal, attempts, goals, primary and secondary assists and individual expected goals per 60 minutes; average shot distance; hits, hits taken, blocked shots, giveaways, takeaways, penalties taken and penalties drawn per 60 minutes; faceoff percentage (shrunken) and faceoffs per 60 minutes; the share of on-ice faceoffs in the offensive zone; height and weight; position indicators (centre, left wing, right wing, defence); and handedness. Birth country from the roster records defines national eligibility.

\paragraph{Reference rosters and opponents.} The basketball references are the twelve-player rosters announced by USA Basketball and Canada Basketball for Paris 2024. Canada's roster included two players with no NBA minutes in 2023--24 (Khem Birch and Melvin Ejim), who are represented at replacement level. The opponents follow the original study: the United States is evaluated against Canada, France and the 2023--24 Oklahoma City Thunder, and Canada against the United States, France and Oklahoma City. France is represented by its twelve NBA players with the most 2023--24 minutes, since most of its Olympic roster played outside the NBA. The hockey reference is Canada's 4 Nations roster, including injury replacements; because a hockey team dresses twelve forwards and six defencemen, we score the twelve forwards and six defencemen with the most 2024--25 time on ice before the tournament. Canada is evaluated against the actual United States, Sweden and Finland rosters, selected in the same way.
